
\begin{frame}[fragile]{La normalisation}

La normalisation, c'est l'art de créer des schémas relationnels où toutes
les relations sont en troisième forme normale, et sans perte d'information.

\vfill

Dans cette session:

\begin{redbullet}
  \item Exemple de décomposition d'un schéma
  \item Algorithme de normalisation
  \item Approche globale
\end{redbullet}


\vfill
\begin{blocgauche}
\alert{Ces diapositives correspondent au support en ligne disponible sur
le site \url{http://sql.bdpedia.fr/}}
\end{blocgauche}

\end{frame}


\begin{frame}{Point de départ: relation globale et dépendances}

On part d'un schéma contenant tous les attributs connus.

\vfill

$$Appart(idAppart, surface, idImmeuble, nbEtages, dateConstruction)$$

\vfill

On identifie les dépendances fonctionnelles

$$idAppart \to  surface, idImmeuble, nbEtages, dateConstruction$$

et

$$idImmeuble \to nbEtages, dateConstruction$$

\vfill
\begin{blocgauche}
\begin{remblock}{En troisième forme normale?}
Non, car la seconde DF montre
une dépendance  dont la partie gauche n'est pas la clé, 
\texttt{idAppart}. \end{remblock}
\end{blocgauche}
\end{frame}

%%%%%%%
\begin{frame}{La décomposition}


On identifie les dépendances fonctionnelles \alert{minimales} et \alert{directes}.

$$idAppart \to  surface, idImmeuble$$

et

$$idImmeuble \to nbEtages, dateConstruction$$

\vfill

On crée une relation pour chacune:

\begin{redbullet}
  \item Appart(\textbf{idAppart}, surface, idImmeuble)
  \item Immeuble (\textbf{idImmeuble}, nbEtages, dateConstruction)
\end{redbullet}


\begin{blocgauche}
On obtient des relations en 3FN, \alert{sans perte d'information}.
\end{blocgauche}
\end{frame}

%%%%%%%%%%%%
\begin{frame}{Regardons les occupants}

Le schéma global de départ est le suivant:


$$Occupant(idPersonne, nom, pr\acute{e}nom, idAppart, surface)$$

\vfill

On a les dépendances suivantes :
$$idAppart \to surface$$
et
$$idPersonne \to pr\acute{e}nom, nom$$

\vfill

La clé est la paire $(idPersonne, idAppart)$. 

\vfill
\begin{blocgauche}
\begin{remblock}{En troisième forme normale?}
Non, car pour les \alert{deux} DF,  la partie gauche n'est pas la clé
\end{remblock}
\end{blocgauche}
\end{frame}

%%%%%%%%%%%%
\begin{frame}{La décomposition}

Comme avant, à partir des DF minimales et directes.

\vfill


\begin{redbullet}
  \item Personne(\textbf{idPersonne}, prénom, nom)
  \item Appart (\textbf{idAppart}, surface)
\end{redbullet}

\vfill

\alert{Pas suffisant}, car on a perdu le lien entre les appartements et les personnes.

\vfill
On ajoute une relation avec la clé.

\begin{redbullet}
  \item Occupant (\textbf{idPersonne, idAppart})
\end{redbullet}

\vfill
\begin{blocgauche}
\alert{On est en 3FN, sans perte d'information (jointures pour reconstituer)}
\end{blocgauche}
\end{frame}

%%%%%%%%%%%%
\begin{frame}{Algorithme de normalisation}

On part d’un schéma de relation $R$ \alert{global}
et d’un ensemble de dépendances fonctionnelles minimales et directes.

\vfill

On détermine alors les clés de R

\vfill 
\begin{redbullet}
  \item Pour chaque DF minimale et directe $X \to A_1 \cdots, A_n$,
      on crée une relation $(X,A_1 \cdots, A_n)$ de clé $X$
  \item Pour chaque clé $C$
non représentée dans une des relations précédentes, on crée une relation $(C)$ de clé $C$.
\end{redbullet}

\vfill

\begin{blocgauche}
\alert{On obtient un schéma normalisé}
\end{blocgauche}
\end{frame}


%%%%%%%%%%%%
\begin{frame}{Et en pratique ?}

Pas tout à fait suffisant: les identifiants n'existent pas naturellement dans la vraie vie...

\vfill

$$(titre, ann\acute{e}e, pr\acute{e}nomMES, nomMES, ann\acute{e}eNaiss)$$

\vfill 

Pas de DF... Il faut les \alert{ajouter} et décidant des \alert{entités} et 
de leur identifiant.

\vfill

Ici, entités Film et Réalisateur, avec \texttt{idFilm} 
et \texttt{idRéalisateur}. Soit: 


$$(idFilm, titre, ann\acute{e}e, idR\acute{e}alisateur, pr\acute{e}nomMES, nomMES, ann\acute{e}eNaiss)$$

avec 

$$idR\acute{e}alisateur \to pr\acute{e}nomMES, nomMES, ann\acute{e}eNaiss$$
et
$$idFilm \to titre, ann\acute{e}e, idR\acute{e}alisateur$$

\begin{blocgauche}
Maintenant on normalise et on obtient un schéma en 3FN.
\end{blocgauche}
\end{frame}


\begin{frame}{Illustration: table de départ}

À partir de cette table pleine d'anomalies.

\begin{small}
\begin{center}
\begin{tabular}{l|l|l|c|l|l|c}
 idFilm & titre &  année &  idRéalisateur & nomMES &  prénomMES &  annéeNaiss \\
\hline
1 & Alien & 1979 & 101 & Scott & Ridley & 1943\\
2 & Vertigo  & 1958 & 102 & Hitchcock & Alfred & 1899\\
3 & Psychose & 1960 &102 &  Hitchcock & Alfred & 1899\\
4 & Kagemusha & 1980 & 103 & Kurosawa & Akira & 1910\\
5 & Volte-face & 1997 & 104 & Woo & John & 1946\\
6 & Pulp Fiction & 1995 & 105 & Tarantino & Quentin & 1963 \\
7 & Titanic & 1997 & 106 & Cameron & James & 1954\\
8 & Sacrifice & 1986 & 107 & Tarkovski & Andrei &1932\\
\hline
\end{tabular}
\end{center}
\end{small}

\end{frame}

\begin{frame}{On obtient après normalisation}


\begin{center}
\begin{minipage}[b]{6cm}
\begin{center}
\begin{tabular}{l|l|c|c}
\textbf{idFilm} & titre & année  & \textit{idR} \\
\hline
1 &Alien & 1979 & 101 \\
2 &Vertigo  & 1958 & 102\\
3 &Psychose & 1960 & 102 \\
4 &Kagemusha & 1980 & 103\\
5 &Volte-face & 1997 & 104 \\
6&Pulp Fiction & 1995 &  105 \\
7&Titanic & 1997 & 106 \\
8&Sacrifice & 1986 & 107\\
\hline
\end{tabular}

La table des films
\end{center}
\end{minipage} \ \hspace*{0.3cm}
\begin{minipage}[b]{5cm}
\begin{center}
\begin{tabular}{l|l|l|c}
{\bf idR} &  nom & prénom &  annéeNaiss \\
\hline
101 & Scott & Ridley & 1943\\
102 & Hitchcock & Alfred & 1899\\
103 & Kurosawa & Akira & 1910\\
104 & Woo & John & 1946\\
105 & Tarantino & Quentin & 1963 \\
106 & Cameron & James & 1954\\
107 & Tarkovski & Andrei &1932\\
\hline
\end{tabular}

La table des réalisateurs
\end{center}
\end{minipage}
\end{center}

\vfill

\begin{blocgauche}
\alert{Schéma normalisé et sans perte d'information}
\end{blocgauche}
\end{frame}


%%%%%%%%%%%%%
\begin{frame}{À retenir}

Il est \alert{toujours} possible de se ramener à un schéma normalisé.

\vfill
\begin{redbullet}
   \item En déterminant les DF et les clés
   \item En appliquant l'algorithme de normalisation
\end{redbullet}

\vfill

\alert{En pratique}: les DF ne sont pas données naturellement. Il faut les déterminer
par un processus de conception basé sur les entités et leurs identifiants.

\vfill
\begin{blocgauche}
\alert{Démarche globale}: conception entité / association,
suivie de la normalisation.
\end{blocgauche}
\end{frame}

